% =====================================================================
% IV. System Design  (틀 — 내용 새로 작성)
% =====================================================================
\section{System Design}\label{sec:design}

% ---------------------------------------------------------------------
\subsection{Recovery-Aware DP ($\psi$+$R$ joint optimization)}\label{sec:design:dp}
% [가이드] 핵심 최적화 정식화. rank 목적함수, DP 점화식, 메모리 제약, 알고리즘.
\todo{문제 정식화 서술 (기호 정의: $\psi$ 배치, $R$ 복구 라우팅 등).}

\begin{equation}\label{eq:rank}
  \text{TODO: rank / objective 정의}
\end{equation}

\begin{equation}\label{eq:dp}
  \text{TODO: DP recurrence } h(j,s)
\end{equation}

\begin{equation}\label{eq:mem}
  \text{TODO: backup-memory reservation 제약}
\end{equation}

\begin{algorithm}[t]
  \caption{Alternating $\psi$+$R$ optimization}
  \label{alg:alt}
  \begin{algorithmic}[1]
    \State \todo{알고리즘 의사코드}
  \end{algorithmic}
\end{algorithm}

% ---------------------------------------------------------------------
\subsection{Mirror cache and chain-aware recovery}\label{sec:design:mirror}
% [가이드] 비동기 미러 캐시 + 백업 승격 시 KV 상태 재구성.
\todo{미러 캐시 / 복구 메커니즘 서술.}

% ---------------------------------------------------------------------
\subsection{Asynchronous chain forwarding}\label{sec:design:async}
% [가이드] ResultReady 역채널, 동기 unwind 제거.
\todo{비동기 체인 포워딩 서술.}

% ---------------------------------------------------------------------
\subsection{Runtime and implementation}\label{sec:design:runtime}
% [가이드] 구현 스택, 코디네이터/워커 역할, 배포 방식.
\todo{런타임/구현 서술.}
